% ============================================================
% Chapter 6: Experimental Design
%   Section: Construction of the Custom Validation Set
%     Subsection: Label Vocabulary
% Included from chapters/06-experimental-design/03-custom-validation-set/custom-validation-set.tex
% ============================================================

\subsection{Label Vocabulary}
\label{subsec:custom-labels}

X-Fact preserves each fact-checker's original label
(\texttt{labelRaw}) next to a normalised one (\texttt{label}); the
custom set does the same. The annotation guide names its labels after
AVeriTeC \cite{schlichtkrull2023averitec}, the real-world,
open-web claim verification benchmark whose four classes ---
\emph{Supported}, \emph{Refuted}, \emph{Not Enough Evidence} and
\emph{Conflicting Evidence/Cherry-picking} --- are defined by the
evidence available rather than by an absolute notion of truth, which
is the stance this work also takes. A fifth label,
\emph{partially supported}, is added because the system's own
\texttt{Verdict} distinguishes a claim whose central assertion holds
while a detail does not, the outcome in which real verification most
often lands. Table~\ref{tab:custom-labels} gives the correspondence.

\begin{table}[ht]
\centering
\small
\begin{tabularx}{\linewidth}{@{}llX@{}}
\toprule
\textbf{\texttt{labelRaw}} & \textbf{\texttt{label}} &
\textbf{Definition} \\
\midrule
supported & \texttt{TRUE} & The evidence supports the claim as
stated. \\
partially supported & \texttt{PARTIALLY\_TRUE} & The central claim
holds; a detail does not. \\
conflicting evidence/cherrypicking & \texttt{MISLEADING} & Sources
of equal standing conflict, or the claim is true but selected to
mislead. \\
refuted & \texttt{FALSE} & The evidence refutes the claim. \\
not enough evidence & \texttt{UNVERIFIED} & There is not enough
independent evidence to decide either way. \\
\bottomrule
\end{tabularx}
\caption{Annotation labels, after AVeriTeC
\cite{schlichtkrull2023averitec}, and the system verdicts they map
to. The X-Fact mapping used in this work sends X-Fact's
\emph{partly true/misleading} to \texttt{MISLEADING} and its
\emph{complicated/hard to categorise} to \texttt{UNVERIFIED}, so both
sets are scored on the same five classes.}
\label{tab:custom-labels}
\end{table}
